%\documentclass[12pt,letterpaper]{article}

%%%
%% Required packages:
%%
\usepackage{etex}
\usepackage{amsmath}
\usepackage{amssymb}
\usepackage{amstext}
\usepackage{amsthm}
\usepackage{fancyhdr,fancybox}
\usepackage{mathrsfs} % need for \mathscr command
\usepackage{multicol}
\usepackage[normalem]{ulem}
%\usepackage{pictex,pictexplus}
% \usepackage{pictexwd}
\usepackage{rotating}
\usepackage{url}
\usepackage{xfrac}
\usepackage{booktabs}
\usepackage{color,soul}
\usepackage{comment}

\usepackage{hyperref}
\hypersetup{
    colorlinks=true,     % false: boxed links; true: colored links
    linkcolor=blue,      % color of internal links
    citecolor=blue,      % color of links to bibliography
    filecolor=blue,      % color of file links
    urlcolor=blue        % color of external links
}


\usepackage{tikz}
\usetikzlibrary{backgrounds,calc}

%%
%% Common formatting commands
%%
\setlength{\parindent}{0in}
\setlength{\textwidth}{7in}
\setlength{\evensidemargin}{-0.25in}
\setlength{\oddsidemargin}{-0.25in}
\setlength{\parskip}{.5\baselineskip}
\setlength{\topmargin}{-0.5in}
\setlength{\textheight}{9in}

%               Problem and Part
\newcounter{problemnumber}
\newcounter{partnumber}
\newcounter{subpartnumber}
\newcommand{\Problem}{%
  \refstepcounter{problemnumber}%
  \setcounter{partnumber}{0}%
  \setcounter{subpartnumber}{0}%
  \item[\makebox{\hfil\textbf{\theproblemnumber.}\hfil}]%
  }
\newcommand{\InProblem}[2][1.6]{%
  \refstepcounter{problemnumber}%
  \setcounter{partnumber}{0}%
  \setcounter{subpartnumber}{0}%
  \textbf{\theproblemnumber.}\ \ \parbox[t]{#1in}{#2}%
}
\newcommand{\InSmallProblem}[1]{\InProblem[1.05]{#1}}
\newcommand{\NoProblem}[1]{%
  \mbox{\hfil\textbf{#1}\hfil}%
  }
\newcommand{\UnProblem}{%
  \refstepcounter{problemnumber}%
  \setcounter{partnumber}{0}%
  \setcounter{subpartnumber}{0}%
  \mbox{\hfil\textbf{\theproblemnumber.}\hfil}%
  }
\newcommand{\InFlexProblem}[1]{%
  \refstepcounter{problemnumber}%
  \setcounter{partnumber}{0}%
  \setcounter{subpartnumber}{0}%
  \textbf{\theproblemnumber.}\ \ #1%
}
%%  Part:
\newcommand{\Part}{%
  \renewcommand{\thepartnumber}{(\alph{partnumber})}%
  \refstepcounter{partnumber}
  \setcounter{subpartnumber}{0}%
  \item[(\alph{partnumber})]%
}
\newcommand{\InPart}[2][1.75]{%
  \renewcommand{\thepartnumber}{(\alph{partnumber})}%
  \refstepcounter{partnumber}%
  \setcounter{subpartnumber}{0}%
  (\alph{partnumber})\ \ \parbox[t]{#1in}{#2}%
}
\newcommand{\UnPart}{%
  \refstepcounter{partnumber}%
  \renewcommand{\thepartnumber}{(\alph{partnumber})}%
  \setcounter{subpartnumber}{0}%
  (\alph{partnumber})%
}
\newcommand{\InLargePart}[1]{\InPart[3.05]{#1}}
\newcommand{\InFullPart}[1]{\InPart[6.2]{#1}}
\newcommand{\InSmallPart}[1]{\InPart[1.05]{#1}}
%% SubPart:
\newcommand{\SubPart}{%
  \refstepcounter{subpartnumber}%
  \item[(\roman{subpartnumber})]%
  }
\newcommand{\InSubPart}[2][2.25]{%
  \stepcounter{subpartnumber}%
  (\roman{subpartnumber})\ \ \parbox[t]{#1in}{#2}%
}
\newcommand{\InSmallSubPart}[1]{\InSubPart[1.5]{#1}}
\newcommand{\InTinySubPart}[1]{%
  \stepcounter{subpartnumber}%
  (\roman{subpartnumber})\ \ \makebox{#1}%
}
\newcommand{\InMySubPart}[2][2.25]{%
  \stepcounter{subpartnumber}%
  \parbox[t]{#1in}{\makebox[0.3in][l]{(\roman{subpartnumber})}\ \ {#2}}%
}
\newcommand{\TrueFalse}{\textbf{T}\ \ \textbf{F}\ \ }
\newcommand{\TFPart}[1]{% 
  \stepcounter{partnumber}%
  \setcounter{subpartnumber}{0}%
  \item[(\alph{partnumber})] \TrueFalse\parbox[t]{5.5in}{#1}}

\newcommand{\Points}[1]{(#1 points) \ }
\newcommand{\Point}[1]{(#1 point) \ }


%% For Answers:
\newcounter{answernumber}
\newcommand{\Answer}{%
  \refstepcounter{answernumber}%
  \setcounter{partnumber}{0}%
  \setcounter{subpartnumber}{0}%
  \item[\makebox{\hfil\textbf{\theanswernumber.}\hfil}]%
  }


% Setting up the headers for the first page
%\chead{} % will default to what is typed in the file.
\fancyhead[L]{\textsc{Math 22a}}
\fancyhead[R]{\textsc{Fall 2024}}
\fancyfoot[R]{
	\vspace*{\fill}\footnotesize\textcopyright{} \emph{Harvard Math 22a}	
}
\renewcommand{\headrulewidth}{0pt}

%\fancyhead[C]{}
%	\chead{}	
%	\lhead{}
%	\rhead{}
%	\cfoot{}


% Putting copyright on the footer of each page
%\fancypagestyle{secondpage}{
%	\fancyhead[C]{TESTing again} % change this to be blank
%		%\chead{}	
%	\fancyhead[L]{bears} % change this to be blank
%	\fancyhead[R]{dogs} % change this to be blank
%	\fancyfoot[R]{ %keeps the footer the same
%		\vspace*{\fill}\footnotesize\textcopyright{} \emph{Harvard Math 22a}	
%	}
%}
%\renewcommand{\headrulewidth}{0pt}
%\pagestyle{fancy}
\pagestyle{empty}


%\lhead{}
%\rhead{}
%\rfoot{\vspace*{\fill}\footnotesize\textcopyright{} \emph{Harvard Math 22a}}
%\renewcommand{\headrulewidth}{0pt}
%\pagestyle{fancy}




%%
%% Common shortcut commands:
%%
\newcommand{\ds}{\displaystyle}
\newcommand{\sgn}{\operatorname{sgn}}
\renewcommand{\Pr}{\operatorname{P}}
\newcommand{\CPr}[3][{}]{\Pr_{\text{#1}}\left(#2\,|\,#3\right)}

\newcommand{\C}{\mathbb{C}}
\newcommand{\R}{\mathbb{R}}
\newcommand{\Z}{\mathbb{Z}}
\newcommand{\N}{\mathbb{N}}
\newcommand{\Q}{\mathbb{Q}}
\newcommand{\D}{\mathbb{D}}

\newcommand{\rref}{\operatorname{rref}}
\newcommand{\im}{\operatorname{im}}
\newcommand{\rank}{\operatorname{rank}}
\newcommand{\diag}{\operatorname{diag}}
\newcommand{\Span}{\operatorname{span}}
%\renewcommand{\vec}[1]{\mathbf{#1}}
\newcommand{\charpoly}{\mathscr{P}}
\newcommand{\nicefrac}[2]{\sfrac{#1}{#2}} % using xfrac package
\newcommand{\topology}{\mathcal{T}}
\newcommand{\Inn}{\operatorname{Inn}}

\newcommand{\blankbox}[2]{\fbox{\rule{#1}{0in}\rule{0in}{#2}}}

\newenvironment{myindentpar}[1]%
 {\begin{list}{}%
         {\setlength{\leftmargin}{#1}
          \setlength{\rightmargin}{\leftmargin}}%
         \item[]%
 }
 {\end{list}}
\newtheorem{theorem}{Theorem}
\newtheorem*{NStheorem}{Theorem}
\newtheorem{cor}[theorem]{Corollary}
\newtheorem*{claim}{Claim}
\newtheorem{defn}{Definition}

\newenvironment{amatrix}[1]{%
  \left(\begin{array}{@{}*{#1}{c}|c@{}}
}{%
  \end{array}\right)
}

%--------grstep
% For denoting a Gauss' reduction step.
% Use as: \grstep{\rho_1+\rho_3} or \grstep[2\rho_5 \\ 3\rho_6]{\rho_1+\rho_3}
\newcommand{\grstep}[2][\relax]{%
   \ensuremath{\mathrel{
       {\mathop{\longrightarrow}\limits^{#2\mathstrut}_{
                                     \begin{subarray}{l} #1 \end{subarray}}}}}}
\newcommand{\swap}{\leftrightarrow}



\chead{\Large\textbf{Functions, \fbox\S 12}}

%\begin{document}
\thispagestyle{fancy}

\begin{enumerate}
  \Problem
  Which of the following functions are one-to-one (injective); which are onto (surjective)?

  \begin{enumerate}
  \item $f:\R\to\R$ given by $f(x) = x$\vfill
  \item $g:\R\to\R_{\ge0}$ given by $g(x) = x^2$\vfill
  \item $h:\R\smallsetminus\{0\}\to\R$ given by $h(x) = 1+\frac{1}{x}$\vfill
  \item $r:\R\to\R$ given by $r(x) = 3x^3-2$\vfill
  \end{enumerate}

  \Problem Consider the function $f:\Z^2 \to\Z^2$ given by $f(m,n) = (m+n,m+2n)$.
  \begin{enumerate}
  \item Show that $f$ is an injective (that is, one-to-one) function. \vfill\vfill\vfill
  \item Show that $f$ is a surjective (that is, onto) function. \vfill\vfill\vfill
  \end{enumerate}

  \newpage

  \Problem A function $f:S\to T$ is called a \emph{constant function}
  if there exists a $t\in T$ such that for all $s\in S$ we have $f(s) = t$.
  \begin{enumerate} 
    \Part Is the following definition equivalent? Why or why not?
    \begin{quote}
      A function is called a \emph{constant function} if for all $s\in
      S$ there exists a $t\in T$ such that $f(s)=t$.
    \end{quote}
    \textbf{Hint:}\ It isn't.
    \vspace{.7in}
  \end{enumerate}

  For the rest of this problem, let $g:S\to T$ and $h:T\to S$ be
  functions, and assume their composition $h\circ g$ is a constant
  function.
  \begin{enumerate}
    \Part Write out a well-quantified statement to explain what
    ``assume their composition $h\circ g$ is a constant function''
    means.  
    \vspace{.7in}

    \Part  What can you say about $h$ given that $g$ is onto?  Express
    your claim as an if-then statement, and prove it. 
    \vfill
    \vfill

    \Part What can you say about $g$ given that $h$ is one-to-one?
    Express your claim as an if-then statement, and prove it. 
    \vfill
    \vfill
    
    \Part Is it necessarily true that $g\circ h$ is a constant
    function? 
    \vfill
  \end{enumerate}
  

\end{enumerate}

%\end{document}
\begin{comment}
\newpage

\chead{\Large\textbf{Functions -- Answers / Solutions}}
\lhead{}
\rhead{}
\thispagestyle{fancy}

\begin{enumerate}
	\Answer
	\begin{enumerate}
		\Part This is a bijection.  We prove this in two steps, by showing
		first that $f$ is injective and next that $f$ is surjective. 
		
		To show that $f$ is injective, we show that if $f(x_1) = f(x_2)$,
		then $x_1 = x_2$.  But $f(x_1) = x_1$ and $f(x_2) = x_2$, so
		clearly these two statements are the same.
		
		To show that $f$ is surjective, we show that if $y\in\R$, then
		there is an $x\in\R$ with $f(x) = y$.  This is simple: take $x=y$,
		so $f(y) = y$.
		
		\Part The function $g$ is \emph{not}\ injective, since $g(-a) =
		g(a) = a^2$.  But $g$ is surjective: If $y\in\R_{\ge0}$ (that is,
		if $y\in\R$ and $y\ge0$), then $y = g(\sqrt{y})$.  
		
		\Part The function $h$ is injective but not surjective.  
		
		To see injectivity, we assume that $h(x_1) = h(x_2)$ and prove
		that $x_1 = x_2$.  If $h(x_1) = h(x_2)$, then $1+\frac{1}{x_1} =
		1+\frac{1}{x_2}$, from which we find that $\frac{1}{x_1} =
		\frac{1}{x_2}$ and so $x_1=x_2$.  
		
		We can \emph{try}\ to prove surjectivity, and in failing we'll see
		why $h$ is not surjective.  Given $y\in\R$, we want to find $x$ so
		that $h(x) = y$.  Solving $1+\frac{1}{x}=y$ for $x$, we get $x =
		\frac{1}{y-1}$.  This isn't solvable when $y=1$, and staring at
		$h(x) = 1+\frac{1}{x}$ we can realize that $h(x)$ will never equal
		$1$.  Thus $h$ is not surjective.
	\end{enumerate}
	
	
	\Answer 
	\begin{enumerate}
		\Part Suppose $f(x,y)=f(a,b)$, then $(x+y,x+2y)=(a+b,a+2b)$. Thus we have the two equations $$x+y=a+b$$ and $$x+2y=a+2b$$, and subtracting the first from the second we get $y=b$. Finally plugging into the first equation we get $x=a$. Thus $(a,b)=(x,y)$, so $f$ is injective.
		\Part Let $(a,b)\in \mathbb{Z}^2$. We need to find $(x,y)$ such that $f(x,y)=(a,b)$. Take $$(x,y)=(2a-b,b-a),$$ then $$f(x,y)=f(2a-b,b-a)=(2a-b+b-a,2a-b+2(b-a))=(a,b).$$ Thus $f$ is surjective.
	%	\Part Suppose instead you consider $$f:\mathbb{N}^2 \to \mathbb{N}^2$$ given by $f(x,y)=(x+y,x+2y)$. Note $f$ is no longer surjective since for all $(x,y)\in \mathbb{N}^2$, $f(x,y)\neq (1,1)$.
	\end{enumerate}
	
	
	\Answer
	\begin{enumerate}
		\Part Here we are not making a single choice of $t$ such that $f(x)=t$ for all $x\in S$; we are instead allowing different $t$ for each element $x\in S$. This is just saying that $f(x)\in T$ for all $x\in S$, which we already knew from the definition of a function.
		\Part There exists an element $b\in S$ such that for all $x\in S$ we have $h(g(x))=b$.
		\Part {\bf Claim:} If $g$ is surjective, then $h$ is constant.
		\begin{proof}
			Since $h \circ g$ is a constant function, there exists an element $b\in S$ so that $(h \circ g)(x)=b$ for all $x\in S$. We want to prove that $h$ is constant; namely we want to show that $h(y)=b$ for all $y\in T$.
			
			Let $y\in T$. Since $g$ is surjective we know there exists $a \in S$ to that $g(a)=y$. Now $$h(y)=h(g(a))=b$$ since $h\circ g$ is constant. Thus $h$ must be a constant function when $g$ is surjective.
		\end{proof}
		\Part {\bf Claim:} If $h$ is injective, then $g$ is constant.
		\begin{proof}
			Let $x,y \in S$, then we want to prove that $g(x)=g(y)$. Since $h \circ g$ is a constant function, we know that $h(g(x))=h(g(y))$. Since $h$ is injective, we have $g(x)=g(y)$. Thus $g$ is a constant function.
		\end{proof}
		\Part No, this is not true. Consider $S=\{a,b,c\}=T$ and $g(a)=a,g(b)=b,g(c)=b$ and $h(a)=a, h(b)=a, h(c)=b$, then $h\circ g$ is constant but $g \circ h$ is not.
	\end{enumerate}
	
\end{enumerate}

\end{comment}
%\end{document}


%%% Local Variables: 
%%% mode: latex
%%% TeX-master: t
%%% End: 
